\documentclass[11pt]{article}
\usepackage[T1]{fontenc}
\usepackage[utf8]{inputenc}
\usepackage{lmodern}
\usepackage{microtype}
\usepackage[a4paper,margin=27mm]{geometry}
\usepackage{amsmath,amssymb,amsthm,mathtools,mathrsfs}
\usepackage{booktabs,longtable,array}
\usepackage{enumitem}
\usepackage{xcolor}
\usepackage{hyperref}
\usepackage[nameinlink,noabbrev]{cleveref}
\usepackage{listings}
\usepackage{fancyhdr}
\usepackage{graphicx}

\hypersetup{
  colorlinks=true,
  linkcolor=blue!55!black,
  citecolor=green!35!black,
  urlcolor=blue!55!black,
  pdftitle={Alaoglu--Erdos: New Progress Beyond the Unified Report},
  pdfauthor={OpenAI GPT-5.6 Pro}
}
\pagestyle{fancy}
\fancyhf{}
\lhead{Alaoglu--Erdos: new progress addendum}
\rhead{\thepage}
\setlength{\headheight}{14pt}
\setlist[itemize]{topsep=3pt,itemsep=2pt,parsep=1pt}
\setlist[enumerate]{topsep=3pt,itemsep=2pt,parsep=1pt}
\setlength{\parindent}{0pt}
\setlength{\parskip}{5pt}

\newtheorem{theorem}{Theorem}[section]
\newtheorem{proposition}[theorem]{Proposition}
\newtheorem{lemma}[theorem]{Lemma}
\newtheorem{corollary}[theorem]{Corollary}
\newtheorem{criterion}[theorem]{Criterion}
\theoremstyle{definition}
\newtheorem{definition}[theorem]{Definition}
\newtheorem{problem}[theorem]{Target problem}
\newtheorem{experiment}[theorem]{Reproducible experiment}
\newtheorem{example}[theorem]{Example}
\theoremstyle{remark}
\newtheorem{remark}[theorem]{Remark}
\newtheorem{warning}[theorem]{Warning}

\newcommand{\Ssmooth}{\mathcal S_{2,3}}
\newcommand{\NN}{\mathbb N}
\newcommand{\ZZ}{\mathbb Z}
\newcommand{\QQ}{\mathbb Q}
\newcommand{\RR}{\mathbb R}
\newcommand{\ord}{\operatorname{ord}}
\newcommand{\lcmop}{\operatorname{lcm}}
\newcommand{\dist}{\operatorname{dist}}
\newcommand{\height}{\operatorname{H}}
\newcommand{\fall}[2]{(#1)_{#2}}
\newcommand{\statusproved}{\textcolor{green!35!black}{\textbf{PROVED}}}
\newcommand{\statusconditional}{\textcolor{orange!70!black}{\textbf{CONDITIONAL}}}
\newcommand{\statusexploratory}{\textcolor{blue!60!black}{\textbf{EXPLORATORY}}}
\newcommand{\statusblocked}{\textcolor{red!60!black}{\textbf{BARRIER}}}

\title{\textbf{Alaoglu--Erdos: New Progress Beyond the Unified Report}\\[3pt]
\large Return-time arithmetic, integral divided differences, jet rigidity, and multi-return alternants}
\author{Research addendum prepared with OpenAI GPT-5.6 Pro}
\date{22 August 2026}

\begin{document}
\maketitle

\begin{center}
\fbox{\begin{minipage}{0.93\textwidth}
\textbf{Status.} This addendum does not claim a completed proof of the Alaoglu--Erdos conjecture. It establishes several rigorous new reductions and obstruction theorems, gives an exact hierarchy of integer-valued exclusion certificates, and isolates a concrete determinant/denominator gap that a successful proof must close. Unverified announcements by other systems are not used as mathematical premises.
\end{minipage}}
\end{center}

\begin{abstract}
Assume that $x\in\RR$ and that $A=2^x$ and $B=3^x$ are integers. The supplied unified report already develops many transcendence-theoretic, analytic, $p$-adic, canonical-product, and approximation-theoretic approaches. The present document is deliberately an addendum rather than a restatement. Its first new tool is an \emph{integral divided-difference sieve}: for any ordered $2,3$-smooth nodes, a minimally cleared divided difference of $t^x$ is a nonzero integer whose sign and size are controlled by a generalized mean-value theorem. This upgrades a three-point curvature observation to arbitrary order and produces explicit forbidden neighborhoods of every integer.

The second new tool is a \emph{jet-rigidity theorem}. If $P\in\QQ[U,V]$ is nonzero, then the order of vanishing of $P(2^t,3^t)$ at a canonical point is exactly the ordinary multiplicity of $P$ at the corresponding point of the algebraic torus. In particular, bidegree $(D_1,D_2)$ gives contact order at most $D_1+D_2$; exact quadratic-contact auxiliary polynomials cannot exist. Combining this with exact denominator lower bounds at return times and Baker-type lower bounds for $\|nx\|$ rules out a broad class of single-return auxiliary-polynomial strategies.

A way around that ceiling is to use many return times simultaneously. We construct a positive exponential alternant whose determinant has quadratic Vandermonde vanishing and whose denominator-cleared value is a nonzero integer. A fully explicit estimate shows that the naive square monomial box is only one factor of order $\sqrt N/\log N$ away from contradiction: the archimedean gain is of size $N^2\log N$, whereas generic denominator clearing costs $N^{5/2}$. This identifies denominator sharing, rather than more Taylor expansion, as the central technical target. We also derive a synchronized two-base return-gap recurrence, a local Sidon-rigidity theorem for additive collisions of $2,3$-smooth numbers, and a formalization plan suitable for Lean.
\end{abstract}

\tableofcontents
\newpage

\section{Scope, source audit, and standards of evidence}

The input archive was parsed before this addendum was written. The decompressed source contains
\texttt{20,483} lines and \texttt{485} section-level headings. These statistics describe the
baseline against which ``new'' is being used here.


A case-insensitive lexical audit of the baseline source gives the following occurrence counts for the principal vocabulary of this addendum. A zero is evidence that the terminology is absent; a nonzero count is not by itself duplication, so every related section below states and proves a quantitatively stronger or structurally different result.

\begin{center}
\small
\begin{tabular}{@{}lr@{}}
\toprule
phrase & occurrences in baseline source\\
\midrule
divided difference & 9 \\
Peano kernel & 0 \\
jet rigidity & 0 \\
alternant & 18 \\
total positivity & 8 \\
local Sidon & 0 \\
return gap & 0 \\
Sturmian & 63 \\
\bottomrule
\end{tabular}
\end{center}

This report assumes the notation and broad background of the supplied unified report, but it does not reproduce its long discussions of the rational case, the Gelfond--Schneider reduction, the Four Exponentials barrier, canonical products, or its previously developed curvature and approximation programs. A short baseline is repeated only when logically necessary for a self-contained proof. Every substantive item below is marked implicitly by its mathematical form:
\begin{itemize}
\item a theorem, lemma, proposition, or corollary is proved in the text;
\item a criterion is a proved implication whose hypothesis remains open;
\item a target problem is a sharply formulated missing statement;
\item an experiment is computational evidence and is not promoted to theorem status.
\end{itemize}

The conjecture is
\[
  2^x,3^x\in\ZZ \quad\Longrightarrow\quad x\in\ZZ.
\]
The report does not use competitive claims, news reports, or private messages as evidence. A proof is counted only if its hypotheses and every inferential step can be audited independently.

\section{Minimal baseline and the smooth semigroup}

Assume throughout that
\[
 A=2^x\in\NN,\qquad B=3^x\in\NN.
\]
Since $2^x$ is a positive integer, either $x=0$ or $x\ge 1$. If $x=p/q\in\QQ$ is in lowest terms and $2^{p/q}$ is an integer, then unique factorization applied to $A^q=2^p$ gives $q\mid p$, hence $q=1$. Thus any counterexample is irrational. Gelfond--Schneider then implies that any counterexample is transcendental.

Set
\[
 \Ssmooth=\{2^a3^b:a,b\in\NN_0\}.
\]
The hypothesis propagates multiplicatively:
\begin{equation}\label{eq:smooth-integrality}
  (2^a3^b)^x=A^aB^b\in\NN\qquad(a,b\ge0).
\end{equation}
This elementary observation is the source of every integer certificate below.


\begin{proposition}[Valuation descent]\label{prop:valuation-descent}
Let $u=\nu_2(A)$ and $v=\nu_3(B)$. For every integer $0\le r\le\min\{u,v\}$,
\[
 2^{x-r}=A/2^r\in\NN,\qquad 3^{x-r}=B/3^r\in\NN.
\]
Thus any counterexample descends to a counterexample, relabeled again as $x$, for which
\begin{equation}\label{eq:primitive-valuation}
 \min\{\nu_2(A),\nu_3(B)\}=0.
\end{equation}
The descended exponent remains positive: for nonintegral $x$, both native valuations are at most $\lfloor x\rfloor$.
\end{proposition}

\begin{proof}
The displayed identities are immediate. If $2^u\mid A=2^x$, then $2^u\le2^x$ and hence $u\le x$; because $u$ is integral and $x$ is not, $u\le\lfloor x\rfloor$, and similarly for $v$. Repeating the descent with $r=\min\{u,v\}$ gives \eqref{eq:primitive-valuation}.
\end{proof}

The number
\[
 \alpha=\log_2 3
\]
is irrational by unique factorization and transcendental by Gelfond--Schneider: if $\alpha$ were algebraic irrational, then $2^\alpha=3$ would be transcendental. This transcendence, rather than mere irrationality, is what makes the jet theorem in \cref{sec:jet} exact.

\section{The integral divided-difference sieve}\label{sec:dd}

\subsection{A master integer certificate}

Let $T=(t_0,\ldots,t_k)$ be distinct positive integers, ordered increasingly, and define
\[
 d_j(T)=\prod_{\substack{0\le i\le k\\i\ne j}}(t_j-t_i),
 \qquad
 L(T)=\lcmop_{0\le j\le k}|d_j(T)|.
\]
The minimally denominator-cleared divided difference is
\begin{equation}\label{eq:DT}
 \mathscr D_T(s)=L(T)\sum_{j=0}^{k}\frac{t_j^s}{d_j(T)}.
\end{equation}
The coefficients $L(T)/d_j(T)$ are integers and have no unnecessary common denominator.

\begin{theorem}[Integral divided-difference sieve]\label{thm:dd-master}
Let $T=(t_0<\cdots<t_k)$ consist of elements of $\Ssmooth$. Under the Alaoglu--Erdos hypothesis, $\mathscr D_T(x)$ is an integer. If $x\notin\{0,1,\ldots,k-1\}$, then it is nonzero and has the sign of
\[
 \fall{x}{k}=x(x-1)\cdots(x-k+1).
\]
Moreover, for some $\xi\in(t_0,t_k)$,
\begin{equation}\label{eq:dd-mvt}
 \mathscr D_T(x)=\frac{L(T)}{k!}\fall{x}{k}\,\xi^{x-k}.
\end{equation}
Consequently
\begin{equation}\label{eq:dd-gap}
 1\le \frac{L(T)}{k!}\,|\fall{x}{k}|\,
 \max\{t_0^{x-k},t_k^{x-k}\}.
\end{equation}
\end{theorem}

\begin{proof}
Equation \eqref{eq:smooth-integrality} makes every $t_j^x$ an integer, while every $L(T)/d_j(T)$ is an integer; hence \eqref{eq:DT} is integral. The sum without $L(T)$ is the ordinary $k$th divided difference $[t_0,\ldots,t_k]f$ of $f(t)=t^x$. The generalized mean-value theorem for divided differences gives
\[
 [t_0,\ldots,t_k]f=\frac{f^{(k)}(\xi)}{k!}
 =\frac{\fall{x}{k}}{k!}\xi^{x-k}.
\]
The derivative has a constant nonzero sign on $(0,\infty)$ unless $x$ is one of $0,\ldots,k-1$. Therefore the integer is nonzero; its absolute value is at least one. The final estimate follows by bounding $\xi^{x-k}$ at an endpoint.
\end{proof}

This theorem is stronger than a curvature test in two ways. It works at arbitrary order, and its integer normalization is canonical enough to be optimized over smooth node sets.

\subsection{Explicit forbidden neighborhoods of every integer}

The next consequence is elementary but useful because it is uniform and machine-checkable.

\begin{corollary}[Endpoint windows]\label{cor:endpoint-windows}
Let $T=(1=t_0<t_1<\cdots<t_k)$ be a smooth node set and put $L=L(T)$. Under the conjectural hypothesis, a nonintegral $x$ cannot lie in either of the intervals
\begin{align}
  k-1 &<x<k-1+\frac1L,\label{eq:right-window}\\
  k-1-\frac{k}{L}&<x<k-1.\label{eq:left-window}
\end{align}
At an endpoint not covered by strict inequality, the same conclusion follows whenever direct substitution in \eqref{eq:DT} gives absolute value strictly below one.
\end{corollary}

\begin{proof}
Write $x=k-1+h$ with $0<h<1$. Since $x-k<0$ and $t_0=1$, \eqref{eq:dd-mvt} gives
\[
 |\mathscr D_T(x)|\le \frac{L}{k!}
 h\prod_{r=1}^{k-1}(r+h)<Lh.
\]
Thus $h<1/L$ makes a nonzero integer have absolute value below one. For the left interval write $x=k-1-h$; then
\[
 \frac{|\fall{x}{k}|}{k!}
 =\frac{h\prod_{r=1}^{k-1}(r-h)}{k!}<\frac{h}{k},
\]
and the same argument applies.
\end{proof}

\begin{example}[The complete strip $1<x<2$]
Take $T=(1,2,3,4)$. Here $L(T)=6$ and
\[
 \mathscr D_T(x)=4^x-3\,3^x+3\,2^x-1.
\]
For $1<x<2$, this is a negative nonzero integer. By the mean-value theorem it equals
\[
 x(x-1)(x-2)\xi^{x-3}
\]
for some $1<\xi<4$. Its absolute value is less than
\[
 x(x-1)(2-x)<1,
\]
because $x(2-x)\le1$ and $0<x-1<1$. This is impossible. Thus the whole strip is excluded by one four-node certificate. The argument is structural; it does not enumerate the possible value $A=3$.
\end{example}

\subsection{Finite search for low-denominator certificates}

For each order $2\le k\le10$, the following table gives the node set with smallest $L(T)$ found by exhaustive search among subsets of the first eighteen $2,3$-smooth numbers, with the node $1$ required. The search is finite and exact; no floating-point arithmetic enters. The last two columns are the rigorous widths from \cref{cor:endpoint-windows} around the integer $k-1$.

\begin{center}
\small
\begin{longtable}{@{}c>{\raggedright\arraybackslash}p{45mm}rcc@{}}
\toprule
order $k$ & nodes $T$ & $L(T)$ & right width $1/L$ & left width $k/L$\\
\midrule
\endfirsthead
\toprule
order $k$ & nodes $T$ & $L(T)$ & right width $1/L$ & left width $k/L$\\
\midrule
\endhead
2 & \(\{1,2,3\}\) & 2 & 0.5 & 1 \\
3 & \(\{1,2,3,4\}\) & 6 & 0.166667 & 0.5 \\
4 & \(\{1,2,3,4,6\}\) & 120 & 0.00833333 & 0.0333333 \\
5 & \(\{1,2,3,4,6,8\}\) & 1680 & $0.595\times 10^{-3}$ & 0.00297619 \\
6 & \(\{1,2,3,4,6,8,9\}\) & 5040 & $0.198\times 10^{-3}$ & 0.00119048 \\
7 & \(\{1,2,3,4,6,8,9,12\}\) & 3991680 & $0.251\times 10^{-6}$ & $0.175\times 10^{-5}$ \\
8 & \(\{1,2,3,4,6,8,9,12,18\}\) & 4071513600 & $0.246\times 10^{-9}$ & $0.196\times 10^{-8}$ \\
9 & \(\{1,2,3,4,6,8,9,12,16,18\}\) & 741015475200 & $0.135\times 10^{-11}$ & $0.121\times 10^{-10}$ \\
10 & \(\{1,2,3,4,6,8,9,12,16,18,24\}\) & 818081084620800 & $0.122\times 10^{-14}$ & $0.122\times 10^{-13}$ \\
\bottomrule
\end{longtable}
\end{center}

For reproducibility, here are the corresponding primitive coefficient vectors for the first six displayed orders. An entry $t:c$ means the term $c\,t^x$.

\begin{center}
\small
\begin{longtable}{@{}c>{\raggedright\arraybackslash}p{125mm}@{}}
\toprule
$k$ & node--coefficient pairs\\
\midrule
\endfirsthead
\toprule
$k$ & node--coefficient pairs\\
\midrule
\endhead
2 & \( 1\!:\!1,\;2\!:\!-2,\;3\!:\!1 \) \\
3 & \( 1\!:\!-1,\;2\!:\!3,\;3\!:\!-3,\;4\!:\!1 \) \\
4 & \( 1\!:\!4,\;2\!:\!-15,\;3\!:\!20,\;4\!:\!-10,\;6\!:\!1 \) \\
5 & \( 1\!:\!-8,\;2\!:\!35,\;3\!:\!-56,\;4\!:\!35,\;6\!:\!-7,\;8\!:\!1 \) \\
6 & \( 1\!:\!3,\;2\!:\!-15,\;3\!:\!28,\;4\!:\!-21,\;6\!:\!7,\;8\!:\!-3,\;9\!:\!1 \) \\
7 & \( 1\!:\!-216,\;2\!:\!1188,\;3\!:\!-2464,\;4\!:\!2079,\;6\!:\!-924,\;8\!:\!594,\;9\!:\!-264,\;12\!:\!7 \) \\
\bottomrule
\end{longtable}
\end{center}

The table is not merely a finite-range computation. It supplies a library of exact integer functionals that can be combined with interval arithmetic or with additional congruence information. The rapid growth of $L(T)$ also exposes the obstruction: the smooth set is too sparse for high-order divided differences to stay small automatically.

\subsection{Peano-kernel generalization}

Divided differences are one instance of a broader framework. Let $t_0<\cdots<t_r$ be smooth nodes and $c_0,\ldots,c_r\in\ZZ$. Suppose
\begin{equation}\label{eq:moment-vanish}
  \sum_{j=0}^{r}c_jt_j^q=0\qquad(0\le q<k).
\end{equation}
Then
\[
  \mathscr L_c(x)=\sum_{j=0}^{r}c_jt_j^x
\]
is an integer under the hypothesis. By the Peano-kernel theorem,
\[
 \mathscr L_c(x)=\int_0^\infty K_c(u)\,
 \frac{\fall{x}{k}}{(k-1)!}u^{x-k}\,du
\]
for an explicit compactly supported piecewise-polynomial kernel $K_c$. If $K_c$ has one sign, then $\mathscr L_c(x)$ has a certified nonzero sign whenever $\fall{x}{k}\ne0$. Therefore a successful high-order certificate problem can be stated as an integer optimization problem:
\begin{problem}[Sign-regular smooth-node certificate]\label{prob:peano}
For each strip $(m,m+1)$, find smooth nodes and primitive integer coefficients satisfying \eqref{eq:moment-vanish}, with a one-sign Peano kernel and with
\[
 \sup_{m<s<m+1}|\mathscr L_c(s)|<1.
\]
Such a certificate excludes the entire strip.
\end{problem}
The divided-difference construction is the canonical totally positive solution with $r=k$, but allowing $r>k$ creates genuine optimization freedom.

\section{Local Sidon rigidity for smooth additive collisions}\label{sec:sidon}

The supplied report's low-order curvature idea can be upgraded to a complete quantitative statement for every balanced four-point collision.

\begin{theorem}[Local Sidon rigidity]\label{thm:local-sidon}
Assume $x>1$ and the Alaoglu--Erdos hypothesis. Let
\[
 p,\ p+h,\ q,\ q+h\in\Ssmooth,
 \qquad 0<p<p+h\le q<q+h.
\]
Then
\begin{equation}\label{eq:sidon-defect}
 \Delta_x(p,q;h)=p^x+(q+h)^x-(p+h)^x-q^x
\end{equation}
is a positive integer. More precisely,
\begin{equation}\label{eq:double-integral}
 \Delta_x(p,q;h)=x(x-1)
 \int_0^h\int_p^q (u+v)^{x-2}\,du\,dv.
\end{equation}
If $1<x\le2$, then
\begin{align}
 x(x-1)h(q-p)(q+h)^{x-2}
 &\le \Delta_x(p,q;h)\nonumber\\
 &\le x(x-1)h(q-p)p^{x-2},\label{eq:sidon-bounds-low}
\end{align}
and for $x\ge2$ the endpoint powers in the lower and upper bounds are reversed.
\end{theorem}

\begin{proof}
Every term in \eqref{eq:sidon-defect} is an integer by \eqref{eq:smooth-integrality}. Writing $f(t)=t^x$,
\begin{align*}
 \Delta_x(p,q;h)
 &=\bigl(f(q+h)-f(q)\bigr)-\bigl(f(p+h)-f(p)\bigr)\\
 &=\int_0^h\bigl(f'(q+v)-f'(p+v)\bigr)\,dv\\
 &=\int_0^h\int_p^q f''(u+v)\,du\,dv.
\end{align*}
Since $f''(t)=x(x-1)t^{x-2}>0$, the defect is positive and therefore at least one. Monotonicity of $t^{x-2}$ gives the two-sided bounds.
\end{proof}

Thus the map $s\mapsto s^x$ is strictly Sidon on every equal-sum fiber of ordered smooth pairs: equal input sums with unequal spreads can never remain equal after applying the map. The integer lower bound makes this quantitative.

\begin{corollary}[Microscopic-collision obstruction]
For $1<x<2$, no balanced smooth collision as in \cref{thm:local-sidon} can satisfy
\[
 x(x-1)h(q-p)p^{x-2}<1.
\]
Equivalently, every such collision obeys
\[
 h(q-p)\ge \frac{p^{2-x}}{x(x-1)}.
\]
\end{corollary}

This looks like a route to contradiction: search for increasingly local additive collisions among smooth numbers. The next proposition explains why fixed four-term collisions cannot deliver them.

\begin{proposition}[Projective finiteness barrier]\label{prop:sunit-finite}
Up to common multiplication by an element of the group $\langle2,3\rangle=\{2^a3^b:a,b\in\ZZ\}$ (and, after choosing a primitive integral representative, up to smooth integral scaling), there are only finitely many nondegenerate solutions of
\[
 u_1+u_4=u_2+u_3,
 \qquad u_i\in\Ssmooth,
\]
where no proper subsum vanishes and the two unordered pairs are distinct. Consequently, scaling a fixed collision cannot make the normalized defect in \eqref{eq:sidon-defect} small: scaling by $q\in\Ssmooth$ multiplies the defect by $q^x$.
\end{proposition}

\begin{proof}
Divide the equation by one term. The resulting normalized equation is an $S$-unit equation with three variables in the finitely generated multiplicative group generated by $2$ and $3$. The $S$-unit theorem gives only finitely many nondegenerate normalized solutions. For the second assertion, homogeneity gives
\[
 \Delta_x(qp,qq;qh)=q^x\Delta_x(p,q;h).
\]
Both factors on the right are positive integers.
\end{proof}

\begin{warning}
Any proof based only on exact four-term smooth additive collisions must therefore introduce something genuinely new: support size tending to infinity, approximate rather than exact collisions, varying integer weights, or a normalization with additional divisibility. Merely finding more scaled copies of a known collision cannot work.
\end{warning}

\section{Synchronized return-time arithmetic}\label{sec:return}

Write
\[
 x=m+\theta,\qquad m=\lfloor x\rfloor,\qquad 0<\theta<1
\]
for a hypothetical nonintegral solution. For $n\ge1$ let
\[
 M_n=\lfloor nx\rfloor=nm+\lfloor n\theta\rfloor,
 \qquad \delta_n=\{n\theta\}\in(0,1).
\]
Irrationality of $\theta$ ensures that $\delta_n$ is never zero. Define the two integer return gaps
\begin{equation}\label{eq:return-gaps}
 a_n=A^n-2^{M_n},\qquad b_n=B^n-3^{M_n}.
\end{equation}
They are positive, and
\begin{equation}\label{eq:return-normalized}
 \frac{a_n}{2^{M_n}}=2^{\delta_n}-1,
 \qquad
 \frac{b_n}{3^{M_n}}=3^{\delta_n}-1.
\end{equation}

\begin{proposition}[Exact return-gap calculus]\label{prop:return-calculus}
For every $n\ge1$:
\begin{enumerate}
\item
\begin{equation}\label{eq:gap-lower}
 \delta_n\ge
 \max\left\{
 \log_2(1+2^{-M_n}),
 \log_3(1+3^{-M_n})
 \right\}.
\end{equation}
\item With
\[
 R(t)=\frac{3^t-1}{2^t-1},\qquad R(0)=\log_2 3=\alpha,
\]
one has the exact rational identity
\begin{equation}\label{eq:R-rational}
 \rho_n:=\frac{b_n2^{M_n}}{a_n3^{M_n}}=R(\delta_n)\in\QQ.
\end{equation}
The function $R$ is strictly increasing on $(0,\infty)$; in particular
\[
 \alpha<\rho_n<2.
\]
\item If $c_{n,r}=\lfloor\delta_n+\delta_r\rfloor\in\{0,1\}$, then
\begin{align}
 a_{n+r}={}&2^{M_n}a_r+2^{M_r}a_n+a_na_r
 -(2^{c_{n,r}}-1)2^{M_n+M_r},\label{eq:a-recurrence}\\
 b_{n+r}={}&3^{M_n}b_r+3^{M_r}b_n+b_nb_r
 -(3^{c_{n,r}}-1)3^{M_n+M_r}.
 \label{eq:b-recurrence}
\end{align}
The same binary carry $c_{n,r}$ drives both recurrences, but its correction has coefficient $1$ in base $2$ and coefficient $2$ in base $3$.
\end{enumerate}
\end{proposition}

\begin{proof}
Since $a_n,b_n$ are positive integers, \eqref{eq:return-normalized} gives \eqref{eq:gap-lower}. Dividing the two equations in \eqref{eq:return-normalized} gives \eqref{eq:R-rational}. To see that $R$ is increasing, put $a=\log3>b=\log2$ and compute
\[
 \frac{R'(t)}{R(t)}=
 \frac{a}{1-e^{-at}}-\frac{b}{1-e^{-bt}}.
\]
For fixed $t>0$, the function $z\mapsto z/(1-e^{-zt})$ is strictly increasing because $e^{zt}>1+zt$. The endpoint values are obtained by l'Hopital's rule and direct substitution at $t=1$. Finally, multiply
\[
 A^n=2^{M_n}+a_n,\qquad A^r=2^{M_r}+a_r
\]
and subtract $2^{M_{n+r}}=2^{M_n+M_r+c_{n,r}}$; the base-$3$ identity is identical.
\end{proof}


\begin{proposition}[Exact native-prime height of a return]\label{prop:return-height}
Put $u=\nu_2(A)$ and $v=\nu_3(B)$. For all sufficiently large $n$ one has $M_n>un$ and $M_n>vn$, and
\begin{align}
 \nu_2(a_n)&=un, &
 \nu_3(b_n)&=vn.\label{eq:native-gap-valuations}
\end{align}
Consequently the normalized return coordinates are in lowest native-prime form:
\[
 2^{\delta_n}=\frac{(A/2^u)^n}{2^{M_n-un}},
 \qquad
 3^{\delta_n}=\frac{(B/3^v)^n}{3^{M_n-vn}}.
\]
If $\rho_n=P_n/Q_n$ is reduced with $P_n,Q_n>0$, then
\begin{equation}\label{eq:rho-height-lower}
 P_n\ge2^{M_n-un},\qquad Q_n\ge3^{M_n-vn}.
\end{equation}
In particular both the numerator and denominator of $\rho_n$ grow exponentially with $n$.
\end{proposition}

\begin{proof}
Because $u,v<x$, the inequalities $M_n>un,vn$ hold for all large $n$. The two terms in $A^n-2^{M_n}$ then have unequal $2$-adic valuations $un$ and $M_n$, so the valuation of their difference is the smaller one; this proves $\nu_2(a_n)=un$. The base-$3$ statement is identical. In
\[
 \rho_n=\frac{b_n2^{M_n}}{a_n3^{M_n}},
\]
the numerator has $2$-adic valuation at least $M_n$, while the denominator has exactly $un$ from $a_n$; after reduction, the numerator therefore retains at least $M_n-un$ powers of $2$. Similarly the denominator retains at least $M_n-vn$ powers of $3$.
\end{proof}

The Taylor expansion at the synchronized return is
\begin{equation}\label{eq:R-taylor}
 R(t)=\alpha\left(
 1+\frac{\log3-\log2}{2}t
 +\frac{(\log3-\log2)(2\log3-\log2)}{12}t^2
 +O(t^3)
 \right).
\end{equation}
Thus every good return of $n\theta$ to zero produces a rational approximation $\rho_n$ to $\alpha$ from above. On any fixed interval $0<t\le t_0<1$, strict positivity and continuity of $R'$ give explicit constants $c_0,c_1>0$ such that
\begin{equation}\label{eq:R-comparable}
 c_0\delta_n\le \rho_n-\alpha\le c_1\delta_n.
\end{equation}

The synchronized recurrences \eqref{eq:a-recurrence}--\eqref{eq:b-recurrence} are a more promising object than the normalized identities alone: they mix archimedean smallness, exact integrality, and a common Sturmian carry sequence. A proof would have to exploit divisibility or gcd information lost after dividing by $2^{M_n}$ and $3^{M_n}$.

\section{Jet rigidity and the single-return barrier}\label{sec:jet}

\subsection{Exact contact cannot be superlinear}

The most important new negative result in this addendum is that exact high-order auxiliary-polynomial contact is impossible for a simple structural reason.

\begin{theorem}[Jet rigidity]\label{thm:jet-rigidity}
Let $P\in\QQ[U,V]$ be nonzero and define
\[
 F_P(t)=P(2^t,3^t).
\]
Then
\begin{equation}\label{eq:jet-equality}
 \ord_{t=0}F_P(t)=\operatorname{mult}_{(1,1)}P,
\end{equation}
where the right side is the least total degree of a nonzero term in $P(1+X,1+Y)$. In particular, if $P$ has bidegree at most $(D_1,D_2)$, then
\begin{equation}\label{eq:contact-ceiling}
 \ord_{t=0}F_P(t)\le D_1+D_2.
\end{equation}
The same statement holds at every integer parameter $t=m$. It also holds at a hypothetical counterexample parameter $t=x$, because $(2^x,3^x)=(A,B)\in\QQ^2$.
\end{theorem}

\begin{proof}
Write
\[
 P(1+X,1+Y)=H_d(X,Y)+H_{d+1}(X,Y)+\cdots,
\]
where $H_d\ne0$ is homogeneous of least total degree $d$. Since
\[
 2^t-1=(\log2)t+O(t^2),\qquad
 3^t-1=(\log3)t+O(t^2),
\]
the coefficient of $t^d$ in $F_P(t)$ is
\[
 H_d(\log2,\log3)
 =(\log2)^d H_d\!\left(1,\log_2 3\right).
\]
The polynomial $H_d(1,T)\in\QQ[T]$ is nonzero. Since $\alpha=\log_2 3$ is transcendental, this coefficient cannot vanish. Hence the order is exactly $d$. The bidegree bound gives $d\le D_1+D_2$.

At $t=m$, replace $P(U,V)$ by $P(2^mU,3^mV)$, which still has rational coefficients, and translate the parameter. At a counterexample $t=x$, use $P(AU,BV)$ instead.
\end{proof}

\begin{corollary}[No exact quadratic jet]
A square bidegree-$D$ auxiliary polynomial cannot have contact order $\asymp D^2$ with the exponential curve at a rational point on that curve. Its exact contact order is at most $2D$, attained for example by $(U-1)^D(V-1)^D$ at $(1,1)$.
\end{corollary}

This explains why a naive dimension count with $(D+1)^2$ coefficients is misleading. The tangent slope is transcendental, so rational coefficients cannot cancel the first nonzero homogeneous component along that slope. The only exact vanishing comes from ordinary algebraic multiplicity at the point.

\subsection{Exact denominator lower bound at a return}

At a return time, put
\[
 u_n=2^{\delta_n}=\frac{A^n}{2^{M_n}},
 \qquad
 v_n=3^{\delta_n}=\frac{B^n}{3^{M_n}}.
\]

\begin{lemma}[Single-return denominator]\label{lem:single-denom}
Let $P\in\ZZ[U,V]$ have bidegree at most $(D_1,D_2)$. Then
\[
 2^{M_nD_1}3^{M_nD_2}P(u_n,v_n)\in\ZZ.
\]
Consequently, if $P(u_n,v_n)\ne0$, then
\begin{equation}\label{eq:denom-lower}
 |P(u_n,v_n)|\ge 2^{-M_nD_1}3^{-M_nD_2}.
\end{equation}
\end{lemma}

\begin{proof}
Clear the largest possible power of $2$ and $3$ in each monomial. The resulting expression is an integer, and a nonzero integer has absolute value at least one.
\end{proof}

If $P$ has multiplicity $r$ at $(1,1)$ and coefficient height $H$, elementary Taylor bounds on $0<\delta\le1/2$ give
\begin{equation}\label{eq:jet-upper}
 |P(2^\delta,3^\delta)|
 \le C^{D_1+D_2}(D_1+1)(D_2+1)H\,\delta^r
\end{equation}
for an absolute effective $C>1$. Combining \eqref{eq:denom-lower} and \eqref{eq:jet-upper}, a proof that forces $P(u_n,v_n)=0$ by size alone would need
\begin{equation}\label{eq:single-return-needed}
 r\log\frac1{\delta_n}
 >M_n(D_1\log2+D_2\log3)
 +O(D_1+D_2+\log H).
\end{equation}
Jet rigidity gives $r\le D_1+D_2$, so the left side has at most linear degree gain, exactly like the denominator cost.

\subsection{Baker lower bounds make the mismatch decisive}

For a hypothetical counterexample, $A/2^m$ is not a rational power of $2$; otherwise $A$ would be a power of $2$ and $x$ would be rational. Hence $A/2^m$ and $2$ are multiplicatively independent algebraic numbers. A standard lower bound for a nonzero linear form in two logarithms yields effective constants $c(A)>0$ and $\kappa(A)>0$ such that
\begin{equation}\label{eq:baker-return}
 \|n\theta\|\ge c(A)n^{-\kappa(A)}
 \qquad(n\ge1).
\end{equation}
In particular
\[
 \log\frac1{\delta_n}=O_A(\log n)
\]
for returns to zero, whereas $M_n\asymp_x n$. Therefore \eqref{eq:single-return-needed} cannot hold for all sufficiently large $n$ for any family whose coefficient height is subexponential in $n(D_1+D_2)$.

\begin{proposition}[Single-return exact-jet barrier]\label{prop:single-jet-barrier}
A strategy that uses one return point $(u_n,v_n)$, an integer polynomial auxiliary function, exact vanishing at $(1,1)$, the generic denominator lower bound \eqref{eq:denom-lower}, and only coefficient-height estimates cannot close the conjecture. Increasing the bidegree does not repair the exponent balance, because exact contact and denominator cost are both linear in the degree while return quality is only polynomial in $n$.
\end{proposition}

\begin{proof}
This is the comparison of \eqref{eq:contact-ceiling}, \eqref{eq:single-return-needed}, and \eqref{eq:baker-return}. Dividing the required inequality by $D_1+D_2$ leaves a left side $O_A(\log n)$ and a right side bounded below by a positive multiple of $M_n\asymp n$, apart from coefficient-height terms that only make the requirement harder.
\end{proof}

\begin{remark}
The proposition is a barrier theorem for a method, not a proof of the conjecture. It identifies the permitted escape routes: unexpected denominator cancellation, approximate rather than exact jets with arithmetically tiny leading coefficients, several return points giving determinant-level vanishing, or an independent nonarchimedean input.
\end{remark}

\section{A multi-return exponential alternant}\label{sec:alternant}

Several return points can evade the linear jet ceiling because determinants acquire a Vandermonde factor of quadratic order.

Let
\[
 \Lambda=\{(a_i,b_i):0\le i<N\}\subset\NN_0^2,
 \qquad
 \lambda_i=a_i\log2+b_i\log3,
\]
with the $\lambda_i$ ordered increasingly. Distinct exponent pairs give distinct $\lambda_i$. Choose distinct return times $n_j$ and order their fractional parts
\[
 0<\delta_0<\delta_1<\cdots<\delta_{N-1}\le\varepsilon<1,
 \qquad M_j=\lfloor n_jx\rfloor.
\]
Consider
\begin{equation}\label{eq:alternant}
 \mathcal A=\det\bigl(e^{\lambda_i\delta_j}\bigr)_{0\le i,j<N}
 =\det\bigl(2^{a_i\delta_j}3^{b_i\delta_j}\bigr)_{i,j}.
\end{equation}

\begin{theorem}[Integral positive alternant]\label{thm:alternant}
Suppose $0\le a_i\le D_1$ and $0\le b_i\le D_2$. Under the Alaoglu--Erdos hypothesis,
\begin{equation}\label{eq:alternant-integer}
 \mathcal I=
 \left(\prod_{j=0}^{N-1}2^{D_1M_j}3^{D_2M_j}\right)\mathcal A
\end{equation}
is a positive integer. If $\lambda_{N-1}\le L$ and $\delta_{N-1}\le\varepsilon$, then
\begin{equation}\label{eq:alternant-upper}
 0<\mathcal A\le
 e^{NL\varepsilon}N^{N/2}
 \frac{(L\varepsilon)^{N(N-1)/2}}
 {\prod_{r=0}^{N-1}r!}.
\end{equation}
Consequently
\begin{align}
 1\le{}&
 2^{D_1\sum_jM_j}3^{D_2\sum_jM_j}
 e^{NL\varepsilon}N^{N/2}
 \frac{(L\varepsilon)^{N(N-1)/2}}
 {\prod_{r=0}^{N-1}r!}.
 \label{eq:alternant-master-ineq}
\end{align}
\end{theorem}

\begin{proof}
Each matrix entry is
\[
 e^{\lambda_i\delta_j}
 =\frac{A^{a_in_j}B^{b_in_j}}
 {2^{a_iM_j}3^{b_iM_j}}.
\]
Multiplying column $j$ by $2^{D_1M_j}3^{D_2M_j}$ makes every entry integral, proving integrality of \eqref{eq:alternant-integer}. Using \cref{prop:return-height}, this can be sharpened after native-prime cancellation to exponents $D_1(M_j-u n_j)$ and $D_2(M_j-v n_j)$; the refinement changes constants but not the asymptotic exponent discussed below. The kernel $e^{uv}$ is strictly totally positive on ordered positive arguments; equivalently, the functions $e^{\lambda_i t}$ form an extended complete Chebyshev system. Hence $\mathcal A>0$.

For the upper bound, transform the columns from function values at $\delta_0,\ldots,\delta_{N-1}$ to Newton divided differences. The determinant picks up the Vandermonde factor
\[
 V(\delta)=\prod_{i<j}(\delta_j-\delta_i),
 \qquad 0<V(\delta)\le\varepsilon^{N(N-1)/2}.
\]
The entry in divided-difference column $r$ is a divided difference of $e^{\lambda_i t}$ of order $r$, and therefore has absolute value at most
\[
 \frac{\lambda_i^r e^{\lambda_i\varepsilon}}{r!}
 \le \frac{L^r e^{L\varepsilon}}{r!}.
\]
Hadamard's inequality bounds the transformed determinant by the product of the Euclidean column norms, giving \eqref{eq:alternant-upper}. Multiplying by the clearing factor and using that $\mathcal I$ is a positive integer gives \eqref{eq:alternant-master-ineq}.
\end{proof}

\subsection{The square-root gap}

Take the square monomial box $0\le a,b\le D$. Then
\[
 N=(D+1)^2,\qquad L\le D\log6.
\]
Even without clustering, all fractional parts lie in an interval of length one. Taking the first $N$ return times gives
\[
 \sum_jM_j=O_x(N^2).
\]
The logarithm of the denominator-clearing factor in \eqref{eq:alternant-master-ineq} is therefore
\[
 O_x(DN^2)=O_x(N^{5/2}).
\]
On the other hand, the superfactorial satisfies
\[
 \log\prod_{r=0}^{N-1}r!
 =\frac12N^2\log N-\frac34N^2+O(N\log N),
\]
so the archimedean alternant gain is of order $N^2\log N$. The method misses contradiction by a factor of order
\[
 \frac{\sqrt N}{\log N}.
\]

This is substantially closer than the single-return method: the determinant has recovered quadratic vanishing. But generic column-by-column denominator clearing introduces one extra factor of the side length $D\asymp\sqrt N$.

\begin{criterion}[Denominator-sharing breakthrough]\label{crit:denom-sharing}
A normalization of the multi-return alternant that saves essentially the full side-length factor $D\asymp\sqrt N$ in the cost
\[
 (D_1+D_2)\sum_jM_j,
\]
while preserving a nonzero integer quotient and the Vandermonde gain, would cross the present asymptotic barrier. More precisely, a cost $o(N^2\log N)$, or an $O(N^2\log N)$ cost with a sufficiently small leading constant, is sufficient in the square-box regime; an $O(N^2\operatorname{polylog}N)$ target is the natural first milestone.
\end{criterion}

The criterion is concrete: one need not improve Taylor expansion or find exponentially good returns. One must prove a shared-divisor statement for the determinant, or choose a basis in which the common denominator is dramatically smaller than the product of the column denominators.

\subsection{Why crude clustering does not repair the balance}

Pigeonholing $K$ fractional parts into an interval of length $\varepsilon$ can provide $N$ points only when $K\varepsilon\gtrsim N$. Their indices and hence their $M_j$ then grow like $K$. The gain $N^2\log(1/\varepsilon)$ is offset by a denominator cost at least of order $\sqrt N\,N K$. Substituting $K\gtrsim N/\varepsilon$ makes the latter polynomial in $1/\varepsilon$, while the former is only logarithmic. Thus ordinary pigeonhole clustering makes the denominator problem worse. The missing ingredient must be arithmetic reuse of denominators, not denser recurrence alone.

\section{Gcd structure and rational approximants}\label{sec:gcd}

The rational approximation in \eqref{eq:R-rational} has reduced denominator
\begin{equation}\label{eq:rho-denominator}
 Q_n=
 \frac{a_n3^{M_n}}
 {\gcd(a_n3^{M_n},\,b_n2^{M_n})}.
\end{equation}
In fact \cref{prop:return-height} already gives the unconditional exponential lower bounds $P_n\ge2^{M_n-un}$ and $Q_n\ge3^{M_n-vn}$ for the reduced fraction $P_n/Q_n$. A stronger moving-target gcd theorem would refine the remaining cofactor structure. Nevertheless, \eqref{eq:R-comparable} gives only an error comparable to $\delta_n$, and Baker's theorem makes the best guaranteed returns polynomially small. A finite irrationality measure for $\alpha$ gives a lower bound of shape $Q_n^{-\mu}$, which is exponentially smaller than a polynomial return error. Therefore even exact native-prime height, and a moving-target gcd theorem by itself, do not prove the conjecture.

This observation prevents a common false start: ``large denominator plus irrationality of $\alpha$'' is not enough. One also needs high-order approximation with a shared denominator.

\begin{problem}[Moving-target gcd estimate]\label{prob:moving-gcd}
Prove a uniform subexponential bound, outside explicitly classified algebraic subgroups, for
\[
 \gcd\bigl(A^n-2^{M_n},\,B^n-3^{M_n}\bigr),
 \qquad M_n=\lfloor nx\rfloor.
\]
The fixed-exponent theorem of Bugeaud--Corvaja--Zannier for $\gcd(a^n-1,b^n-1)$ suggests the right scale, but the moving target $M_n$ and the coupled two-base structure require an additional argument.
\end{problem}

\begin{criterion}[High-order shared-denominator extrapolation]
Suppose one can build rational numbers $p_{n,D}/q_{n,D}$ from a bounded number of synchronized gap data with
\[
 \left|\alpha-\frac{p_{n,D}}{q_{n,D}}\right|
 \le \delta_n^{cD^2},
 \qquad
 \log q_{n,D}\le CnD+O(D^2),
\]
and with no catastrophic cancellation that invalidates a lower bound for the reduced denominator. Choosing $D\gg n/\log n$ would then beat a finite irrationality measure for $\alpha$. Jet rigidity shows that such quadratic order cannot come from one exact polynomial jet; a determinant or a genuinely approximate arithmetic jet is required.
\end{criterion}

\section{Two-orbit and carry formulations}\label{sec:dynamics}

Two exact reformulations may be useful for future work.

\subsection{A common carry word}

Set
\[
 X_n=2^{\delta_n}=\frac{A^n}{2^{M_n}},
 \qquad
 Y_n=3^{\delta_n}=\frac{B^n}{3^{M_n}}.
\]
Let
\[
 c_n=\lfloor(n+1)\theta\rfloor-\lfloor n\theta\rfloor\in\{0,1\}.
\]
Then
\begin{equation}\label{eq:carry-dynamics}
 X_{n+1}=\frac{X_nX_1}{2^{c_n}},
 \qquad
 Y_{n+1}=\frac{Y_nY_1}{3^{c_n}}.
\end{equation}
For irrational $\theta$, $(c_n)$ is a Sturmian mechanical word. Thus a counterexample is equivalent to two rational multiplicative codings, in bases $2$ and $3$, driven by exactly the same aperiodic carry word. This is not a proof, but it suggests a Cobham-type rigidity question in which recognizability would have to come from divisibility of the numerators $A^n$ and $B^n$, not merely from symbolic dynamics.

\subsection{Two lattice-point orbits on one transcendental curve}

The curve
\[
 y=t^\alpha,\qquad \alpha=\log_2 3,
\]
contains the canonical integer orbit $(2^n,3^n)$ for $n\ge0$. A counterexample would add the independent integer orbit $(A^n,B^n)$. Their parameters are $n$ and $nx$. Diophantine approximation makes these parameters recurrently close, but only polynomially close; the corresponding affine lattice points remain exponentially far apart. Any lattice-geometry proof must therefore use high-order determinants or projective normalization, not Euclidean proximity of individual points.

\section{What has been ruled out, and what remains viable}\label{sec:scorecard}

\begin{center}
\small
\begin{tabular}{@{}p{36mm}p{38mm}p{72mm}@{}}
\toprule
Route & Status & Precise conclusion\\
\midrule
Fixed three- or four-point curvature & \statusproved & Extends to the arbitrary-order divided-difference sieve and full local Sidon rigidity.\\
Exact four-term smooth collisions & \statusblocked & Finitely many projective shapes; scaling multiplies the defect by $q^x$ and cannot create smallness.\\
Single-return exact auxiliary polynomial & \statusblocked & Jet order is at most total degree, while generic denominator cost is linear in degree times $M_n\asymp n$.\\
Rational approximation from one gap ratio & \statusblocked & Error is $\asymp\delta_n$, much larger than the exponential scale associated with its denominator.\\
Multi-return alternant & \statusexploratory & Restores quadratic Vandermonde vanishing; naive denominator clearing is one factor $\sqrt N/\log N$ too expensive.\\
Moving-target gcd plus shared denominator & \statusconditional & A strong enough quotient-integrality theorem could close the alternant gap.\\
Approximate arithmetic jets & \statusconditional & Must make the leading homogeneous form unusually small at $\alpha$; exact cancellation is forbidden by jet rigidity.\\
\bottomrule
\end{tabular}
\end{center}

The most plausible immediate target is not another analytic estimate. It is an arithmetic divisibility theorem for the alternant or its Vandermonde factorization.

\section{A concrete next-stage research program}\label{sec:program}

\subsection{Target A: factor the denominator-cleared alternant}

For return points in an arithmetic progression $n_j=jq$ before the first carry, the alternant becomes a generalized Vandermonde in the rational quantities
\[
 R_{a,b}=2^{a\delta_q}3^{b\delta_q}
 =\frac{A^{aq}B^{bq}}{2^{aM_q}3^{bM_q}}.
\]
The determinant factors into products of differences $R_{a,b}-R_{a',b'}$. After clearing denominators, each factor has an integer numerator
\[
 A^{aq}B^{bq}2^{a'M_q}3^{b'M_q}
 -A^{a'q}B^{b'q}2^{aM_q}3^{bM_q}.
\]
The open task is to identify a large systematic common divisor of the full product that can be divided out while retaining integrality. A gain of one power of the monomial side length $D$ in the logarithmic denominator cost would be asymptotically decisive.

\subsection{Target B: approximate jets with certified leading forms}

By \cref{thm:jet-rigidity}, exact cancellation of the lowest homogeneous component is impossible unless the component itself vanishes. The alternative is to choose a polynomial whose leading form $H$ has a very small, but nonzero, value
\[
 H(\log2,\log3)=(\log2)^dH(1,\alpha).
\]
This turns the problem into an explicit transcendence-measure optimization for $\alpha$. A useful lemma would jointly control degree, coefficient height, the smallness of $H(1,\alpha)$, and the denominator of the evaluated auxiliary polynomial. Treating these quantities separately loses the needed factor.

\subsection{Target C: moving-target gcd with carry sensitivity}

The recurrences \eqref{eq:a-recurrence}--\eqref{eq:b-recurrence} contain more information than a generic pair of recurrence sequences because the same carry appears in both bases. One should study prime valuations of the pair $(a_n,b_n)$ across additions $n+r$, separating carry and no-carry cases. A theorem showing that a prime cannot divide too many synchronized gap pairs without forcing a multiplicative relation among $2,3,A,B$ would feed directly into \cref{crit:denom-sharing}.

\subsection{Target D: optimized one-sign Peano kernels}

The finite search behind the table in \cref{sec:dd} minimizes only $L(T)$ for pure divided differences. A more flexible integer program should optimize the actual supremum norm on each strip under moment constraints and a one-sign Peano kernel. Total positivity turns the sign constraint into a tractable variation-diminishing condition. Even if this does not prove all strips, it can provide exact finite-range certificates with much smaller support than brute-force enumeration of $A$.

\section{Lean formalization blueprint}\label{sec:lean}

The proved core is suitable for formalization in small independent modules.

\subsection{Module graph}

\begin{enumerate}
\item \texttt{Smooth23.lean}: define $\Ssmooth$, prove closure under multiplication, and prove \eqref{eq:smooth-integrality} using positivity and the real-power product law.
\item \texttt{DividedDifferenceInt.lean}: define $d_j(T)$ and $L(T)$ for a finite ordered list of natural nodes; prove the cleared coefficients are integers and establish integrality of $\mathscr D_T(x)$.
\item \texttt{DividedDifferenceSign.lean}: formalize the generalized mean-value theorem by induction from Rolle's theorem if the exact library lemma is absent; specialize to $t\mapsto t^x$ and prove \cref{thm:dd-master}.
\item \texttt{LocalSidon.lean}: prove the double-integral identity and monotonic endpoint bounds in \cref{thm:local-sidon}.
\item \texttt{ReturnGaps.lean}: formalize $M_n$, $\delta_n$, positivity of $a_n,b_n$, the rational identity \eqref{eq:R-rational}, monotonicity of $R$, and the synchronized recurrences.
\item \texttt{JetRigidity.lean}: represent the translated polynomial by homogeneous components and prove the leading coefficient is nonzero from transcendence of $\log_2 3$. This module needs a theorem packaging the Gelfond--Schneider consequence.
\item \texttt{Alternant.lean}: formalize denominator clearing and positivity via a Chebyshev-system or total-positivity lemma; prove the Hadamard/divided-difference upper bound.
\end{enumerate}

\subsection{Suggested theorem signatures}

\begin{lstlisting}[language={},basicstyle=\ttfamily\small,breaklines=true]
theorem smooth_rpow_integral
    (hx2 : IsInteger (2 ^ x)) (hx3 : IsInteger (3 ^ x))
    {a b : Nat} : IsInteger (((2^a) * (3^b) : Real) ^ x)

 theorem cleared_dividedDifference_int
    (hT : forall t in T, Smooth23 t)
    (hx : AEHypothesis x) : IsInteger (DDclear T x)

 theorem cleared_dividedDifference_sign
    (hord : StrictMono T) (hxk : x notin Finset.range k) :
    sign (DDclear T x) = sign (fallingFactorial x k)

 theorem jet_order_eq_multiplicity
    (P : MvPolynomial (Fin 2) Rat) (hP : P != 0) :
    orderAtZero (fun t => evalCurve P t) = multAtOneOne P
\end{lstlisting}

The formalization should keep the S-unit finiteness proposition and Baker estimates as imported high-level theorems or explicitly labeled assumptions at first. The elementary core---integrality, divided differences, return recurrences, and the denominator-cleared alternant---can be checked independently.

\section{Conclusions}

The work in this addendum changes the shape of the problem in three concrete ways.

First, the integrality of $s^x$ on the smooth semigroup is not limited to a three-term curvature constraint. It yields a complete hierarchy of integer divided differences and one-sign Peano functionals. These give exact exclusion certificates and a principled optimization problem.

Second, exact auxiliary-polynomial contact has a rigid ceiling. Because $\log_2 3$ is transcendental, the exponential curve cannot have anomalously high rational-algebraic contact at a rational point. This removes an entire class of seemingly promising ``many coefficients, many vanishing derivatives'' arguments and explains exactly why they fail.

Third, multi-return alternants recover the missing quadratic order through a Vandermonde determinant. Their present failure is no longer vague: the archimedean determinant gain is $N^2\log N$, while generic denominator clearing costs $N^{5/2}$. A shared-denominator or determinant-divisibility theorem saving one factor of the monomial side length would cross the barrier. That is the most sharply defined route emerging from this round of work.

No final proof is asserted. The report supplies proved lemmas, explicit inequalities, reproducible finite searches, a formalization blueprint, and a narrow arithmetic target rather than an unsupported claim of resolution.

\appendix
\section{Derivation of the Taylor coefficient in the gap ratio}

Put $a=\log3$ and $b=\log2$. Then
\[
 \frac{e^{at}-1}{e^{bt}-1}
 =\frac ab
 \frac{1+\frac a2t+\frac{a^2}{6}t^2+O(t^3)}
 {1+\frac b2t+\frac{b^2}{6}t^2+O(t^3)}.
\]
Using $(1+u)^{-1}=1-u+u^2+O(u^3)$ gives
\[
 \frac{e^{at}-1}{e^{bt}-1}
 =\frac ab\left[
 1+\frac{a-b}{2}t
 +\frac{2a^2+b^2-3ab}{12}t^2+O(t^3)
 \right],
\]
and $2a^2+b^2-3ab=(a-b)(2a-b)$, proving \eqref{eq:R-taylor}.

\section{Exact finite-search specification}

The table in \cref{sec:dd} is produced as follows.
\begin{enumerate}
\item Generate the first eighteen distinct integers of the form $2^a3^b$.
\item For each $2\le k\le10$, enumerate every $(k+1)$-element subset containing $1$.
\item Compute $d_j(T)=\prod_{i\ne j}(t_j-t_i)$ exactly and set $L(T)=\lcmop_j|d_j(T)|$.
\item Minimize lexicographically $(L(T),\max T,\sum T)$.
\item Verify exactly that the primitive coefficients $L(T)/d_j(T)$ annihilate the moments $0,\ldots,k-1$.
\end{enumerate}
The finite pool restriction is stated explicitly; the table makes no claim of global optimality over all smooth nodes.

\section{References}
\begin{thebibliography}{99}

\bibitem{AlaogluErdos1944}
L. Alaoglu and P. Erdos,
\emph{On highly composite and similar numbers},
Transactions of the American Mathematical Society \textbf{56} (1944), 448--469.

\bibitem{Baker1975}
A. Baker,
\emph{Transcendental Number Theory},
Cambridge University Press, 1975.

\bibitem{BakerWustholz2007}
A. Baker and G. Wuestholz,
\emph{Logarithmic Forms and Diophantine Geometry},
Cambridge University Press, 2007.

\bibitem{BombieriGubler2006}
E. Bombieri and W. Gubler,
\emph{Heights in Diophantine Geometry},
Cambridge University Press, 2006.

\bibitem{BCZ2003}
Y. Bugeaud, P. Corvaja, and U. Zannier,
\emph{An upper bound for the G.C.D. of $a^n-1$ and $b^n-1$},
Mathematische Zeitschrift \textbf{243} (2003), 79--84.

\bibitem{deBoor2005}
C. de Boor,
\emph{Divided differences},
Surveys in Approximation Theory \textbf{1} (2005), 46--69.

\bibitem{Evertse1984}
J.-H. Evertse,
\emph{On equations in $S$-units and the Thue--Mahler equation},
Inventiones Mathematicae \textbf{75} (1984), 561--584.

\bibitem{Karlin1968}
S. Karlin,
\emph{Total Positivity, Vol. I},
Stanford University Press, 1968.

\bibitem{Lang1966}
S. Lang,
\emph{Introduction to Transcendental Numbers},
Addison--Wesley, 1966.

\bibitem{LMN1995}
M. Laurent, M. Mignotte, and Y. Nesterenko,
\emph{Formes lineaires en deux logarithmes et determinants d'interpolation},
Journal of Number Theory \textbf{55} (1995), 285--321.

\bibitem{Waldschmidt2000}
M. Waldschmidt,
\emph{Diophantine Approximation on Linear Algebraic Groups},
Springer, 2000.

\end{thebibliography}

\end{document}
